\section{A cutoff polylogarithm identity and its Euler primitive}
\label{hier:gf-section}

The first rank-two member of the new hierarchy gives a compact
polylogarithmic consequence.  Its weight is
\begin{equation}\label{hier:gf-weight}
 R(n)=\frac1{(n+1)^2}-\frac1{(n+2)^2},
\end{equation}
whereas the predecessor's explicit repeated-pole generator uses the
distinct second-difference weight
$\tfrac12\{(n+1)^{-2}-2(n+2)^{-2}+(n+3)^{-2}\}$.
Define the finite ordered determinant sum
\begin{equation}\label{hier:gf-cutoff}
 Z_N=\sum_{0\le n<k\le N}(k-n)^2R(n)R(k),\qquad Z_0=0.
\end{equation}

Our one-direction multiple-polylogarithm convention is
\begin{equation}\label{hier:gf-convention}
 \operatorname{Li}_{s_1,\ldots,s_j}(z)
 =\sum_{n_1>\cdots>n_j\ge1}
       \frac{z^{n_1}}{n_1^{s_1}\cdots n_j^{s_j}},
 \qquad |z|<1,
\end{equation}
with positive integer indices.  Thus only the largest summation index
carries $z$; in a multiple-argument convention this is
$\operatorname{Li}_{s_1,\ldots,s_j}(z,1,\ldots,1)$.
In particular the shuffle identity has normalization
$\operatorname{Li}_1(z)^2=2\operatorname{Li}_{1,1}(z)$.

\begin{theorem}[Finite cutoff and its complete generator]
\label{hier:gf-theorem}
For $M=N+1\ge1$,
\begin{equation}\label{hier:gf-finite}
 \boxed{\displaystyle
 Z_{M-1}=
 2\left(1-\frac1{(M+1)^2}\right)H_M
 -(H_M^{(2)})^2
 -\frac{M^2}{(M+1)^2}\bigl(H_M^{(2)}+1\bigr).}
\end{equation}
The generating function $G(z)=\sum_{N\ge0}Z_Nz^{N+1}$ satisfies
\begin{align}\label{hier:gf-closed}
 z(1-z)G(z)={}&2\operatorname{Li}_1(z)-\operatorname{Li}_2(z)
                 -z-z\operatorname{Li}_4(z)\notag\\
 &+2(1-z)\bigl(\operatorname{Li}_{1,2}(z)
                         -\operatorname{Li}_{2,1}(z)\bigr)
 -(1+z)\operatorname{Li}_{2,2}(z).
\end{align}
All series are absolutely and locally normally convergent for $|z|<1$.
The quotient in \eqref{hier:gf-closed} is interpreted by its removable
value at zero; $G(z)=\tfrac5{48}z^2+O(z^3)$.
\end{theorem}
\begin{proof}
Use the monic moment basis $1,n+1$, which leaves the determinant
unchanged, and put $p=M+1$.  Telescoping gives its three moments as
\[
 \mu_0=1-p^{-2},\qquad
 \mu_1=H_M^{(2)}-Mp^{-2},\qquad
 \mu_2=2H_M-H_M^{(2)}-1+2p^{-1}-p^{-2}.
\]
Expanding $\mu_0\mu_2-\mu_1^2$ proves \eqref{hier:gf-finite}.
The basic absolutely convergent summation identities are
\begin{align*}
 \sum_{M\ge1}H_M^{(p)}z^M
    &=\frac{\operatorname{Li}_p(z)}{1-z},\qquad
 \sum_{M\ge1}\frac{z^M H_M^{(p)}}{(M+1)^q}
    =z^{-1}\operatorname{Li}_{q,p}(z),\\
 \sum_{M\ge1}(H_M^{(2)})^2z^M
    &=\frac{2\operatorname{Li}_{2,2}(z)+\operatorname{Li}_4(z)}{1-z},
 \qquad
 \sum_{M\ge1}\frac{z^M}{(M+1)^q}
    =z^{-1}\operatorname{Li}_q(z)-1.
\end{align*}
The square identity follows by splitting its two inner indices into
the two strict orderings and the diagonal.  Substitute these formulas
into \eqref{hier:gf-finite}, after writing
$M^2/(M+1)^2=1-2/(M+1)+1/(M+1)^2$.  Collecting terms gives
\eqref{hier:gf-closed}.

Equation \eqref{hier:gf-finite} gives $Z_N=O(\log(N+2))$.
Every multiple-polylogarithm series in \eqref{hier:gf-convention}
also converges normally on compact subdisks, since its inner sums
grow at most polynomially while $z^{n_1}$ decays geometrically.
These estimates justify all interchanges.  The original generating
series proves that the apparent singularity at zero is removable;
$Z_0=0$ and $Z_1=R(0)R(1)=5/48$ give the stated initial terms.
\end{proof}

\begin{corollary}[An anchored Euler primitive]
\label{hier:gf-primitive}
For $|z|<1$, set
$I(z)=\int_0^zG(t)\,dt/t$ along any path in the disk.  Then
\begin{align}\label{hier:gf-primitive-closed}
 I(z)={}&\frac{
 \operatorname{Li}_1(z)^2+2\operatorname{Li}_{2,1}(z)
 -\operatorname{Li}_{1,2}(z)+\operatorname{Li}_{2,2}(z)
 +\operatorname{Li}_2(z)-\operatorname{Li}_1(z)}{z}\notag\\
 &-\operatorname{Li}_5(z)-\operatorname{Li}_{1,4}(z)
     -2\operatorname{Li}_{3,2}(z)
     -2\operatorname{Li}_{1,2,2}(z).
\end{align}
The removable value at zero is $I(0)=0$, and
\[
 I(z)=\sum_{N\ge0}\frac{Z_N}{N+1}z^{N+1}
     =\frac5{96}z^2+O(z^3).
\]
\end{corollary}
\begin{proof}
Let $\theta=z\,d/dz$.  From the absolutely convergent series,
\[
 \theta\operatorname{Li}_{p,\boldsymbol s}
      =\operatorname{Li}_{p-1,\boldsymbol s}\quad(p>1),\qquad
 \theta\operatorname{Li}_{1,\boldsymbol s}
      =\frac{z}{1-z}\operatorname{Li}_{\boldsymbol s}.
\]
For an empty $\boldsymbol s$ the second formula means
$\theta\operatorname{Li}_1=z/(1-z)$.
These identities also prove
$\operatorname{Li}_1^2=2\operatorname{Li}_{1,1}$ by comparing
derivatives and values at zero.  Apply them to
\eqref{hier:gf-primitive-closed}: its Euler derivative is the function $G(z)$ specified by
\eqref{hier:gf-closed}.  The convergent expansions at zero give the
zero integration constant and the stated removability.  Alternatively,
coefficientwise division of \eqref{hier:gf-finite} by $M$ gives the
same formula by
\[
 \frac1{M(M+1)^2}
  =\frac1M-\frac1{M+1}-\frac1{(M+1)^2}.
\]
\end{proof}

The raw generator has a representation of depth at most two, although
its coefficient determinant has harmonic degree two.  Its displayed
Euler primitive uses depth at most three.  These are explicit
representations, not claims of minimal depth under functional identities.

\subsection{A convergent second Euler endpoint}

\begin{corollary}[An ordinary-zeta endpoint]
\label{hier:gf-endpoint}
The second Euler primitive has the absolutely convergent value
\begin{align}\label{hier:gf-endpoint-value}
 \sum_{N\ge0}\frac{Z_N}{(N+1)^2}
 &=\int_0^1\frac{-\log z}{z}G(z)\,dz\notag\\
 &=1+3\zeta(2)-2\zeta(3)-\frac34\zeta(4)
                     -\zeta(3)^2-\frac{19}{24}\zeta(6).
\end{align}
The new content here is the determinant endpoint identity; the
weight-six multiple-zeta reductions used below are classical.
\end{corollary}
\begin{proof}
The estimate $Z_N=O(\log(N+2))$ proves convergence.  Positivity of
$Z_N$ also permits Tonelli's theorem in the integral representation.
Divide \eqref{hier:gf-finite} by $M^2$ and collect its rational
coefficients to obtain the pointwise identity
\[
 \frac{Z_{M-1}}{M^2}
 =-\frac{(H_M^{(2)})^2}{M^2}
  -\frac{2H_M}{(M+1)^2}
  +\frac{4H_M}{M(M+1)}
  -\frac{H_M^{(2)}+1}{(M+1)^2}.
\]
Every term is separately summable.  The standard nested-series
interpretation and elementary telescoping give
\[
 \sum_{M\ge1}\frac{H_M}{M(M+1)}=\zeta(2),\qquad
 \sum_{M\ge1}\frac{H_M}{(M+1)^2}=\zeta(2,1)=\zeta(3),
 \qquad
 \sum_{M\ge1}\frac{H_M^{(2)}}{(M+1)^2}=\zeta(2,2)=\frac34\zeta(4).
\]
Splitting coincident and strictly ordered indices gives
\begin{equation}\label{hier:gf-square-Euler}
 \sum_{M\ge1}\frac{(H_M^{(2)})^2}{M^2}
 =2\zeta(2,2,2)+2\zeta(4,2)+\zeta(2,4)+\zeta(6)
 =\zeta(3)^2+\frac{19}{24}\zeta(6).
\end{equation}
For clarity, the three weight-six inputs are
\[
 \zeta(2,2,2)=\frac3{16}\zeta(6),\qquad
 \zeta(4,2)=\zeta(3)^2-\frac43\zeta(6),\qquad
 \zeta(2,4)=\frac{25}{12}\zeta(6)-\zeta(3)^2.
\]
They follow from short classical calculations.  Expanding the cube
of $\zeta(2)$ by coincident indices gives the first.  Euler's formula
$2\zeta(m,1)=m\zeta(m+1)-\sum_{j=1}^{m-2}
\zeta(j+1)\zeta(m-j)$, with $m=5$, gives
$\zeta(5,1)=\tfrac34\zeta(6)-\tfrac12\zeta(3)^2$;
see \cite[Eq.~(1.3)]{hier:BB2005}.  Equating the convergent shuffle
and stuffle products
\[
 \zeta(3)^2=2\zeta(3,3)+6\zeta(4,2)+12\zeta(5,1),
 \qquad
 \zeta(3)^2=2\zeta(3,3)+\zeta(6)
\]
then gives the second weight-six input.  Finally,
$\zeta(2,4)+\zeta(4,2)=\zeta(2)\zeta(4)-\zeta(6)$ gives the third.
The same Euler formula at $m=2$ supplies $\zeta(2,1)=\zeta(3)$.
Substitution into the summable pointwise identity proves
\eqref{hier:gf-endpoint-value}.
\end{proof}
